\scanpage{47}\mspage{34}

ordinal limite. Donc il reste à déterminer les $\unsure{v|u}_{\alpha,\alpha+1}:Z_\alpha\longrightarrow Z_{\alpha+1}$\inserted{ (pour $\alpha<\alpha_0$)}, ce qui se voit sur le diagramme

\diagram{\begin{tikzcd}[ampersand replacement=\&,row sep=large,column sep=large]
 X_{I_\alpha}\arrow[r]\arrow[d] \& X_{I_{\alpha+1}}\arrow[r]\arrow[d] \& X_{I_{\alpha_0}}=Z_0\arrow[d,Rightarrow]\\
 Y_{I_\alpha}\arrow[r,dash] \& Y_{I_\alpha,I_{\alpha+1}}\arrow[r,dash]\arrow[d,"u_{I_\alpha,I_{\alpha+1}}"] \& \deleted{\unsure{Z|Y}}\unsure{Y|Z}_{I_\alpha,I_{\alpha_0}}=Z_\alpha\arrow[d,"v_{\alpha,\alpha+1}"]\\
 \& Y_{I_{\alpha+1}}\arrow[r] \& Y_{I_{\alpha+1},I_{\alpha_0}}=Z_{\alpha+1}.
\end{tikzcd}}

\deleted{\unsure{Donc} Ainsi, \unsure{Rest. Rev.}}

Reprenons la démonstration de 2.1 $\deleted{\unsure{1^\circ|2^\circ}}\unsure{2^\circ|1^\circ}$.

\deleted{Dans le cas actuel $I_0=\varnothing_M$, donc $X_{I_0}\longrightarrow Y_{I_0}$ est $\varnothing_{\unsure{M|0}}\xrightarrow{\sim}\varnothing_M$, donc $Z_0={}$}

On \deleted{\unsure{l'}}a prouvé, par les \inserted{2} lemmes précédents, \deleted{dans le cas particulier $K=\varnothing$, \unsure{On} procède de $=$ dans le cas général\unsure{,} d'un couple $(K,\unsure{J|I})$ quelconque, en utilisant que l'on sait déjà que
\[
 (u_K:X_K\longrightarrow Y_K)\in C.
\]}

\deleted{\unsure{Ceci} fait, en revenant \unsure{au} $I_0=K$, dans $I$.}

Notons que lorsque l'espace sous-jacent à $I$ est noethérien, alors la construction transfinie est forcément finie, donc \unsure{on} n'a pas besoin de la stabilité de $C$ par \unsure{composés|composées} transfinis infinis,

% archive sheet number: 46

